\stepcounter{section}
\section*{Lecture 9: Closed Subschemes after a bit of Sheaf Theory}
Notes by Yuki Ishihara.

Given an abelian group \(G\) and a subgroup \(N \subset G\), we can construct the quotient group \(G/N\). This motivates us to wonder if we can construct the quotient sheaf of abelian groups.  So let $X$ be a topological space, $\mathcal{F}$ a sheaf on $X$ and $\mathcal{G} \subset \mathcal{F}$ a subsheaf of abelian groups. One approach to construct $\mathcal{F}/\mathcal{G}$ is to send the open set $U \to \mathcal{F}(U)/\mathcal{G}(U) =: \mathcal{F}/\mathcal{G}(U)$.  This is a good candidate for the quotient and is a good presheaf however is not a sheaf in general. The following example illustrates how the gluing axiom can fail.
\begin{example}
    $X = \mathbb{P}_\C^1, \mathcal{F} = \mathcal{O}_X$ and $G \subset \mathcal{O}_X$ subsheaf consisting of sections that evaluate to 0 at 0 and $\infty$. Or:
    $$\mathcal{G}(U) = \left\{f \in \mathcal{O}_X : f(0) = 0 \text{ if } 0 \in U \text{ and } f(\infty) = 0 \text{ if } \infty \in U\right\}$$
    We consider the points $0 = [0 : 1] \in \mathbb{P}_\C^1$ and $\infty = [1 : 0] \in \mathbb{P}_\C^1$ and the open set $U = \mathbb{P}^1_\C$. Let $U_1 = U\backslash \{\infty\}$ and $U_2 = U \backslash \{0\}$ an open covering of $U$. \smallskip \\
    On $U$ we have $\mathcal{O}_X(U) = \C$ the constant functions and $\mathcal{G}(U) = \{0\}$ ie. the only constant map that evaluates to 0 at 0 and $\infty$ is the zero map. So $\mathcal{O}_X/\mathcal{G}(U) = \C$ \smallskip \\
    On $U_1 = \spec\C[x]$ we have $\mathcal{O}_X(U_1) = \C[x]$ and $\mathcal{G}(U_1) = x\C[x]$. So $\mathcal{O}_X/\mathcal{G}(U_1) = \C$. \smallskip \\
    On $U_2 = \spec\text{ }\C[x^{-1}]$ we have $\mathcal{O}_X(U_2) = \C[x^{-1}]$ and $\mathcal{G}(U_2) = x^{-1}\C[x^{-1}]$. So $\mathcal{O}_X/\mathcal{G}(U_2) = \C$ \smallskip \\
    On the intersection $U_1 \cap U_2 = \mathbb{P}^1_\C$ we have $\mathcal{O}_X(U_1 \cap U_2) = \C[x, x^{-1}]$ and $\mathcal{G}(U_1 \cap U_2) = \C[x, x^{-1}]$. So $\mathcal{O}_X/\mathcal{G}(U_1 \cap U_2) = 0$. \smallskip \\
    Now take $1 \in \mathcal{O}_X/\mathcal{G}(U_1)$ and $2 \in \mathcal{O}_X/\mathcal{G}(U_2)$. These both restrict to 0 on the intersection trivially. The gluing axiom requires us to find a global section represented by a complex number $f = c$ that restricts to $1$ and $2$ on $U_1$ and $U_2$ respectively. However the restrictions $f|_{U_1} = c$ and $f|_{U_2} = c$ are just the inclusion map and since $1 \neq 2$ there cannot exist such an $f$.
\end{example}
To overcome this issue we use a tool known as sheafification. This follows from the following theorem.
\begin{theorem} (Sheafification) Given a presheaf $\mathcal{P}$ on $X$ there exists a sheaf $\tilde{\mathcal{P}}$ (called the sheafification of $\mathcal{P}$) on $X$ and a map $p : \mathcal{P} \to \tilde{\mathcal{P}}$ such that:
\begin{enumerate}
    \item For any sheaf $\mathcal{Q}$ on $X$ and a map $q : \mathcal{P} \to \mathcal{Q}$ there is a unique map $\tilde{q} : \tilde{\mathcal{P}} \to \mathcal{Q}$ such that the following diagram commutes
        \[\begin{tikzcd}
	{\mathcal{P}} && {\tilde{\mathcal{P}}} \\
	\\
	&& {\mathcal{Q}}
	\arrow["p", from=1-1, to=1-3]
	\arrow["q"', from=1-1, to=3-3]
	\arrow["{\tilde{q}}", from=1-3, to=3-3]
\end{tikzcd}\]
\item For any $x \in X$, the map \(\mathcal{P} \to \tilde{\mathcal{P}}\) induces an isomorphism on stalks $\mathcal{P}_x \cong \tilde{\mathcal{P}}_x$.
\end{enumerate}
\end{theorem}
\newpage 
Considering this we can define:
$$\mathcal{F}/\mathcal{G} := \text{Sheafification of } (U \mapsto \mathcal{F}/\mathcal{G}(U))$$
as the quotient sheaf of abelian groups. It comes along with the following universal property:
\[\begin{tikzcd}
	{\mathcal{G}} && {\mathcal{F}/\mathcal{G}} \\
	& {\mathcal{F}} \\
	& {\mathcal{Q}}
	\arrow["0", from=1-1, to=1-3]
	\arrow["i", from=1-1, to=2-2]
	\arrow["0"', from=1-1, to=3-2]
	\arrow["{\exists!}", from=1-3, to=3-2]
	\arrow["\pi", from=2-2, to=1-3]
	\arrow["q"', from=2-2, to=3-2]
\end{tikzcd}\]
where $i$ is the inclusion map and $\pi$ is the canonical quotient map.
That is, given a sheaf \(\mathcal{Q}\) and a map \(\mathcal{F} \to \mathcal{Q}\) such that the composite \(\mathcal{G} \to \mathcal{Q}\) is zero, there exists a unique map \(\mathcal{F}/ \mathcal{G} \to \mathcal{Q}\) as in the diagram above.

Motivated by group theory we now aim to reproduce the first isomorphism theorem in terms of sheaves. In order to do this we first need to know the following definitions.
\begin{definition} The following are definitions on maps $\mathcal{F} \to \mathcal{Q}$ between sheaves of abelian groups. 
\begin{itemize}
    \item (Surjectivity) The map $\mathcal{F} \to \mathcal{Q}$ is surjective is it is surjective on all stalks. 
    \item (Kernel) The kernel sheaf $\mathcal{K}$ of the map $\mathcal{F} \to \mathcal{Q}$ is the sheaf $\mathcal{K}(U) = \ker(\phi_U)$ where $\phi_U : \mathcal{F}(U) \to \mathcal{Q}(U)$ is the induced map. 
\end{itemize}
\end{definition}
\begin{example}
    The quotient map $\mathcal{F} \to \mathcal{F}/\mathcal{G}$ is surjective sending the stalk $\mathcal{F}_x \mapsto \left[\mathcal{F}_x\right]$. 
\end{example}
We then get the first isomorphism equivalent for sheaves. 
\begin{corollary} (First Isomorphism Theorem)
    If $\mathcal{F} \to \mathcal{Q}$ surjective and $\mathcal{K}$ the kernel sheaf then $\mathcal{F}/\mathcal{K} \to \mathcal{Q}$ is an isomorphism of sheaves. 
\end{corollary}
\begin{exercise}
    Any map on sheaves that is an isomorphism on the stalks is an isomorphism
\end{exercise}
Caution: $\mathcal{F} \to \mathcal{Q}$ surjection of sheaves does not mean $\mathcal{F}(U) \to \mathcal{Q}(U)$ is a surjection for an open set $U$. 
\begin{nonexample}
    $X = \mathbb{P}^1, \mathcal{F} = \mathcal{O}_X$ and $I \subset \mathcal{O}_X$ the subsheaf consisting of sections vanishing at $0$ and $\infty$ then $\mathcal{O}_X \twoheadrightarrow \mathcal{O}_{X/I}$ but $$\C = \mathcal{O}_X(X) \to \mathcal{O}_{X/I}(X) = \C^2$$ is not surjective. 
\end{nonexample}
We now look towards closed embeddings and closed subschemes. Let $X = \spec A$ be an affine scheme and $Y$ a scheme with $i : Y \to X$ a morphism. The following theorem illustrates the main idea given without proof. 
\begin{theorem}
    There exists an ideal $I \lhd A$ with $Y \cong \spec(A/I)$ if and only if the following three conditions hold:
    \begin{enumerate}
        \item The image $i(Y) \subset X$ is closed
        \item $i : Y \to i(Y)$ is a homeomorphism
        \item $i^\sharp : \mathcal{O}_X \to i_*\mathcal{O}_Y$ is surjective
    \end{enumerate}
    In such a case we call $Y$ a closed subscheme of $X$ and $i$ a closed embedding.
\end{theorem}
So in general $i : Y \to X$ is a closed embedding if it is topologically (1 \& 2) and if $i^\sharp$ is surjective (3). The following are two non-examples of this idea. The first not fulfilling (2) and the second not fulfilling (3). 
\begin{nonexample}
    We consider the map $\A^1 \sqcup \A^1 \to \A^2$ where $t$ from the first $\A^1$ gets sent to $(t, 0)$ and $s$ from the second gets sent to $(s, 0)$. This map is not a homeomorphism onto its image as the origin of both $\A^1$ are sent to the origin of $\A^2$ and hence not injective. 
\end{nonexample}
\begin{nonexample}
    (A better non-example) We consider the map $i : \A^1_\C \to \A^2_\C$ where $a \mapsto (a^2, a^3)$. On the level of rings this is equivalent to the map $\C[x, y] \to \C[t]$ where $x \mapsto t^2$ and $y \mapsto t^3$. In this case the image of $i$ is closed and the the map from $\A^1_\C$ to its image is a homeomorphism. Now lets look at the map $i^\sharp : \mathcal{O}_{\A^2_\C} \to i_*\mathcal{O}_{\A^1_\C}$. On the level of stalks, this gives us the map $\C[x, y]_{(x, y)} \to \C[t]_{(t)}$ where $x \mapsto t^2$ and $y \mapsto t^3$. On this stalk however the map is not surjective as $t$ is not in the image. 
\end{nonexample}
We finish of this segment of the notes with the following equivalence stated without proof. 
\begin{proposition}
    TFAE
    \begin{enumerate}
        \item $i : Y \to X$ is a closed embedding
        \item For some open cover $\{U_i\}$ of $X$ we have that $i^{-1}(U_i) \to U_i$ are all closed embeddings
        \item For every open cover $\{U_i\}$ of $X$ we have that $i^{-1}(U_i) \to U_i$ are all closed embeddings
    \end{enumerate}
\end{proposition}
